\chapter{Минимальная архитектура cross-carrier морфизма для c0}
\label{ch:v8-baryon-c0-minimal-cross-carrier-morphism-architecture-gate}

Предыдущий реестр дал $0/7$, но не запрещал новый морфизм. Построим его
минимальный тип. Пусть $L_{\rm src}=\mathbb R u$ --- прямая выбранного
внутреннего инварианта $r$, а $L_{\rm aux}=\mathbb R e$ --- прямая
коэффициента квадратичной массы вспомогательного поля. Калибровочная группа
действует на обеих прямых тривиально:
\begin{equation}
 \rho_{\rm src}(g)u=u,
 \qquad
 \rho_{\rm aux}(g)e=e.
 \label{eq:v8-baryon-c0-trivial-line-actions}
\end{equation}

\section{Классификация морфизмов}

Для двух тривиальных вещественных представлений условие эквивариантности
тождественно, поэтому
\begin{equation}
 \operatorname{Hom}_G(L_{\rm src},L_{\rm aux})\cong\mathbb R,
 \qquad
 \dim\operatorname{Hom}_G=1.
 \label{eq:v8-baryon-c0-equivariant-hom-space}
\end{equation}
Каждый морфизм имеет вид
\begin{equation}
 \mathcal M_\kappa(ru)=\kappa r e,
 \qquad \kappa\in\mathbb R.
 \label{eq:v8-baryon-c0-morphism-family}
\end{equation}
Положительность полюсной массы оставляет $\kappa>0$, но не выбирает точку
этого луча.

\section{Минимальный общий родитель}

После присоединения вспомогательного поля наиболее экономный quadratic
cross-carrier блок равен
\begin{equation}
 S_{\rm cross}[X,\phi]
 =S_{\rm src}[X]
 +\frac12\phi\bigl(z+a\kappa r(X)\bigr)\phi-g\phi J(X).
 \label{eq:v8-baryon-c0-minimal-cross-parent}
\end{equation}
Если источник имеет изолированный вакуум $r(X_*)=r_*>0$, то
\begin{equation}
 \Hess_\phi S_{\rm cross}\big|_{X_*}=z+a\kappa r_*,
 \qquad
 c_0=\kappa r_*.
 \label{eq:v8-baryon-c0-cross-parent-hessian}
\end{equation}
Тем самым архитектура действительно строит требуемую типизированную карту,
но переносит недостающую свободу в её нормировку.

Замена координаты на исходной прямой сохраняет физический коэффициент:
\begin{equation}
 u\mapsto su,qquad r\mapsto r/s,qquad
 \kappa\mapsto s\kappa,qquad
 \kappa r\mapsto\kappa r.
 \label{eq:v8-baryon-c0-source-normalization-orbit}
\end{equation}
Следовательно, число $\kappa$ не имеет смысла без общей нормировки двух
носителей.

\section{Сильнейший кандидат и недостающий селектор}

Для ближайшего суперкривизностного инварианта $r_*=4$ две положительные
эквивариантные карты дают
\begin{equation}
 \kappa=1\Longrightarrow c_0=4,
 \qquad
 \kappa=\frac14\Longrightarrow c_0=1.
 \label{eq:v8-baryon-c0-positive-morphism-witnesses}
\end{equation}
Обе удовлетворяют тем же симметриям и target-independence.

Единственный минимальный способ убрать луч внутри данной архитектуры ---
потребовать изометричность нормированных прямых:
\begin{equation}
 \mathcal M_\kappa^*\mathcal M_\kappa=I_{L_{\rm src}}
 \quad\Longleftrightarrow\quad
 \kappa^2=1
 \quad\Longrightarrow\quad
 \kappa=1\ \hbox{при }\kappa>0.
 \label{eq:v8-baryon-c0-isometric-normalization}
\end{equation}
Но существующий родитель не утверждает, что cross-carrier карта является
изометрией относительно уже выведенных следовых метрик.

\begin{theorem}[Существование и одномерный дефект нормировки]
Минимальный калибровочно-эквивариантный cross-carrier морфизм существует и
единственен с точностью до одного вещественного множителя $\kappa$.
Положительность сокращает пространство до луча $\kappa>0$, но не выбирает
$c_0$. Изометричность выбрала бы $\kappa=1$, однако является новым
условием до явного общего следового вложения.
\end{theorem}

\begin{nogocriterion}[Архитектура не фиксирует нормировку]
Нельзя объявлять каноническую стрелку между двумя тривиальными прямыми
изометрией только потому, что обе одномерны. Требуется вычислить pullback
метрики из одного оператора или действия и проверить
$\mathcal M^*G_{\rm aux}\mathcal M=G_{\rm src}$.
\end{nogocriterion}

\begin{protocol}\begin{enumerate}
\item эквивариантный Hom: \textbf{одномерный};
\item типизированная архитектура: \textbf{существует};
\item положительность фиксирует знак: \textbf{да};
\item положительность фиксирует модуль: \textbf{нет};
\item изометрическая нормировка: \textbf{достаточна, но не выведена};
\item следующий гейт: \textbf{общая следовая нормировка вложения}.
\end{enumerate}\end{protocol}

Машинный сертификат сохранён в
\path{s2t_v8_baryon_c0_minimal_cross_carrier_morphism_architecture_gate_results.json}.